%----------------------------------------------------------------------------
\chapter{Experiments and results}
\label{chap:Experiments}
%----------------------------------------------------------------------------

\begin{itemize}
    \item This chapter reports measurements along the four evaluation axes from \autoref{sec:EvalAxes}: efficiency, robustness, physical realism, flexibility.
    \item All experiments use the shared scaffolding and fairness controls of \autoref{sec:Fairness}: matched gravity, $dt = 1/30$~s, no damping, $\alpha_{\mathrm{XPBD}} = 1/k_{\mathrm{VBD}}$, Chebyshev and adaptive initialisation disabled in VBD for the core comparison.
\end{itemize}

%----------------------------------------------------------------------------
\section{Computational efficiency}
\label{sec:EffResults}
%----------------------------------------------------------------------------

\subsection{Wall-clock}
\label{ssec:ErrorVsMs}
\begin{itemize}
    \item Sweep through different substep and iteration configurations for each solver and plot their Pareto line (optimal solutions) based on avg ms/frame and avg??? Root Mean Squared position error from the Newton reference. Each diagram compares to a Newton with different substep counts
    \item XPBD substep count $n_{\mathrm{sub}} \in \{1, 2, 4, 8, 16, 32\}$
    \item VBD iteration count $n_{\mathrm{it}} \in \{1, 2, 4, 8, 16, 32\}$ 
    \item The simulated problem: a chain with 10 nodes, one end fixed, the other having a weight of ??? (with bending???) . Both ends start from the same height, only gravity acting on it.
    \item VBD reaches the same error value as XPBD with less iterations
    \item VBD's iterations cost more than XPBD's
    \item With this particular problem VBD's first guesses always seem to be closer to reference than XPBD's.
    \item From Newton with 4 substeps graph: VBD with the reference Newton's substeps converges really good. XPBD reaches it's smallest error with relatively few iterations, turning it higher doesn't improve error, only increases compute time. 
    \item From Newton with more substeps: iterations loose their value as with small time steps both solvers converge at 1 or 2 iterations. XPBD takes over because of its cheaper iterations.
\end{itemize}
\emph{TODO: add ss1, ss2, ss4, ss64}

\subsection{Frame time vs.\ node count}
\label{ssec:NodeCountVsFPS}
\begin{itemize}
    \item Measure FPS over a hanging-cloth scene with wind interaction at mesh resolutions \emph{TODO: list resolutions, e.g.\ $8\times 8$ to $128\times 128$}, holding solver settings fixed at the \autoref{sec:Fairness} defaults.
    \item Based on the diagrams in ssec:ErrorVsMs, I choose substep 4 iteration 4 for XPBD and substep 4 iteration 2 for VBD, as these sit at the middle point of "realistic vs. efficient" and are alway close to each other
\end{itemize}
\emph{TODO: Add figure + table}


%----------------------------------------------------------------------------
\section{Robustness and stability}
\label{sec:Robustness}
%----------------------------------------------------------------------------
\subsection{High mass ratio at the end of a pendulum}
\label{ssec:MassRatioPendulum}
\begin{itemize}
    \item A single spring pendulum with a heavy weight at the free end and a light chain leading to the pivot.
    \item Sweep the mass ratio $m_{\mathrm{tip}} / m_{\mathrm{chain}}$ from $1$ to $10^{4}$; report each solver's qualitative behaviour and quantitative position error vs.\ Newton.
    \item This is the mass-disparity regime the primal/dual lens flags for XPBD's per-constraint updates.
\end{itemize}
\emph{TODO: Figure --- error vs.\ mass ratio.}

\subsection{High stiffness ratio chain}
\label{ssec:StiffnessRatioChain}
\begin{itemize}
    \item A chain with alternating stiff and soft springs; sweep the stiffness ratio from $1$ to $10^{4}$.
    \item This is the ill-conditioned-block regime the primal/dual lens flags for VBD's per-vertex update.
\end{itemize}
\emph{TODO: Figure --- error vs.\ stiffness ratio.}

\subsection{Self-collision: cloth folded onto itself}
\label{ssec:SelfCollResults}
\begin{itemize}
    \item A cloth dropped onto a small obstacle, folding onto itself.
    \item Qualitative visual comparison, noting that the contact model differs (projection vs.\ penalty) per \autoref{sec:SelfCollision}.
\end{itemize}
\emph{TODO: Renders, qualitative description.}

%----------------------------------------------------------------------------
\section{Physical realism}
\label{sec:Realism}
%----------------------------------------------------------------------------

\subsection{Newton reference}
\label{ssec:NewtonRef}
\begin{itemize}
    \item Build a high-iteration Newton solver targeting $\|\nabla G\| < 10^{-9}$ (or until further iterations reduce $G$ by less than $10^{-12}$).
    \item The Newton solver shares storage and energies with the production solvers, so the comparison is over algorithm, not material.
    \item Error metric throughout: mass-weighted position error $\|\mathbf{x} - \mathbf{x}^{\star}\|_{\mathbf{M}}$, where $\mathbf{x}^{\star}$ is the Newton solution.
\end{itemize}

\subsection{Harmonic oscillator}
\label{ssec:Harmonic}
\begin{itemize}
    \item A single spring-mass system oscillating in 1D, compared against the analytic solution $x(t) = A\cos(\omega t - \phi)\,e^{-\zeta\omega t}$.
    \item I measure:
    \begin{itemize}
        \item period as a function of $h$ and of iteration / substep count,
        \item amplitude decay rate as a function of damping coefficient.
    \end{itemize}
    \item This is the analogue of Fig.~2 in \cite{10.1145/2994258.2994272}.
\end{itemize}
\emph{TODO: Figure --- measured period and decay, both solvers, against analytic.}

\subsection{Energy tracking}
\label{ssec:EnergyTracking}
\begin{itemize}
    \item For an undamped free-falling and bouncing cloth, plot total mechanical energy over time.
    \item Identify systematic energy gain (instability) or loss (numerical damping) for each solver.
\end{itemize}
\emph{TODO: Figure --- energy vs.\ time, both solvers.}

\subsection{Hanging cloth: stretched area vs.\ reference}
\label{ssec:HangingCloth}
\begin{itemize}
    \item A square cloth pinned at the top corners, allowed to settle.
    \item Measure steady-state stretched area against the Newton reference.
    \item A realism check at static equilibrium, where transient dynamics drop out.
\end{itemize}
\emph{TODO: Figure --- stretched area error vs.\ stiffness.}

\subsection{Damping cross-check}
\label{ssec:DampingCheck}
\begin{itemize}
    \item At matched Rayleigh coefficients $\alpha_{\mathrm{Rayleigh}}$ and $\beta_{\mathrm{Rayleigh}}$, measure the amplitude-decay curve produced by each solver on the harmonic oscillator.
    \item Identical decay confirms the damping models of \autoref{sec:Damping} are calibrated consistently; this both validates fairness and is a result in its own right.
\end{itemize}
\emph{TODO: Figure --- decay envelopes overlaid.}

%----------------------------------------------------------------------------
\section{Flexibility}
\label{sec:FlexResults}
%----------------------------------------------------------------------------

\subsection{Iteration-count independence of stiffness}
\label{ssec:IterIndepStiff}
\begin{itemize}
    \item For each solver, sweep stiffness $k$ over several decades and measure how many iterations / substeps are needed to reach a fixed error threshold.
    \item In the compliance formulation, XPBD's per-iteration behaviour is roughly stiffness-independent (Fig.~6 in \cite{10.1145/2994258.2994272}).
    \item \autoref{fig:IterIndep} shows the settled shape of a pinned cloth at four iteration / substep counts for both solvers. For both methods the resting shape is essentially unchanged as the iteration count grows, so the converged configuration does not depend on how many iterations are spent.
    \item For both solvers the resting configuration is essentially unchanged across the sweep, indicating that the converged shape is independent of the iteration count.
\end{itemize}

\begin{figure}[ht]
    \centering
    \setlength{\tabcolsep}{1.5pt}
    \renewcommand{\arraystretch}{0.4}
    \newcommand{\iterimg}[1]{\raisebox{-0.5\height}{\includegraphics[width=0.215\textwidth,keepaspectratio]{figures/IterationIndependence/#1}}}
    \begin{tabular}{c cccc}
        & {\small $n_{\mathrm{it}}=20$} & {\small $n_{\mathrm{it}}=40$} & {\small $n_{\mathrm{it}}=80$} & {\small $n_{\mathrm{it}}=160$} \\[0.2ex]
        \raisebox{-0.5\height}{\rotatebox[origin=c]{90}{\small VBD}} &
            \iterimg{VBD_i20.png} & \iterimg{VBD_i40.png} & \iterimg{VBD_i80.png} & \iterimg{VBD_i160.png} \\[10ex]
        \raisebox{-0.5\height}{\rotatebox[origin=c]{90}{\small XPBD}} &
            \iterimg{XPBD_i20_16x16.png} & \iterimg{XPBD_i40_16x16.png} & \iterimg{XPBD_i80_16x16.png} & \iterimg{XPBD_i160_16x16.png} \\
    \end{tabular}
    \caption{Settled shape of a pinned cloth at increasing iteration counts ($n_{\mathrm{it}} = 20, 40, 80, 160$).
    Note that in this experiment the material parameters are different for the two solvers.}
    \label{fig:IterIndep}
\end{figure}

\subsection{Stiffness-ratio test}
\label{ssec:StiffRatioTest}
\begin{itemize}
    \item Same kind of plot as the previous subsection, but varying the \emph{ratio} of stiffnesses across the mesh rather than a uniform stiffness.
    \item This is the experiment of \cite{Chen2024} Fig.~24 (and AVBD~\cite{Giles2025} Fig.~4).
    \item Includes a brief discussion of AVBD's improvement on this axis without implementing it myself.
\end{itemize}
\emph{TODO: Figure --- error vs.\ stiffness ratio, with AVBD reference qualitatively cited.}

\subsection{Ease of adding a new constraint or force}
\label{ssec:NewConstraint}
I implement a representative additional force --- \emph{TODO: choose, e.g.\
bending-angle constraint or area-preservation force} --- in both solvers
and report:
\begin{itemize}
    \item lines of code added,
    \item whether a Hessian derivation was required (VBD: yes; XPBD: no, a gradient suffices),
    \item time to first correct simulation.
\end{itemize}
\begin{itemize}
    \item The metric is admittedly subjective but captures a real practitioner cost that the more quantitative experiments above do not.
\end{itemize}
\emph{TODO: Table --- LOC, derivation effort, time-to-correct.}
